% SGA 3, Expose V, Loop 1 English reconstruction.
% Scope: affine/ring setup for Theorem 4.1 and part a) of its proof,
% stopping immediately before part b).
% Authority: Polo--Gille Exp5-13oct24.pdf, component-local p. 12
% (combined-reader p. 274). This re-edited page bears printed header 262
% and internal original-page markers 264--265.
% Jacob Reinhold's CC BY 4.0 Markdown is comparison/drafting material only.

在当前考虑的情形中, \(X_0,X_1,X_2\) 都仿射. 因而可假设
\[
  X_i=\operatorname{Spec}A_i,\qquad
  d_j=\operatorname{Spec}\delta_j,\qquad
  d'_k=\operatorname{Spec}\delta'_k,
\]
其中各 \(A_i\) 为交换环, \(\delta_j,\delta'_k\) 为环同态.
图 \hyperref[sga3:v:lemma:1-2:equation:0-1-2]{(0,1,2)}
可替换为以下对偶图:

% Native Loop-2 reconstruction; local p. 12 / combined p. 274.
\begin{equation*}
  \resizebox{0.57\textwidth}{!}{%
    \SGAThreeNativeInput{figures/exp5/p262_dual_cocartesian_012star.png}}
  \phantomsection\label{fig:sga3-exp5-dual-012star}
\end{equation*}
此图左侧的两个方形为余 Cartesian 方形.

设 \(B\) 是 \(A_0\) 中由满足 \(\delta_0(a)=\delta_1(a)\)
的元素 \(a\in A_0\) 组成的子环.

\medskip
\phantomsection\label{proof:sga3-V-4.1-a}
\noindent\textbf{a)} 先证明 \(A_0\) 在 \(B\) 上整.
对 \(a\in A_0\), 令
\begin{equation}
  P_{\delta_1}(T,\delta_0(a))
  =T^n-\sigma_1T^{n-1}+\cdots+(-1)^n\sigma_n
  \label{eq:sga3-V-4.1-characteristic-polynomial}
\end{equation}
为 \(\delta_0(a)\) 的特征多项式,
其中通过 \(\delta_1\) 将 \(A_1\) 视为 \(A_0\)-代数;
参见 Bourbaki, \emph{Alg\`ebre}, VIII, \S12,
以及 \emph{Alg\`ebre commutative}, II, \S5, 习题~9.
由于图 \hyperref[fig:sga3-exp5-dual-012star]{\((0,1,2)^\ast\)}
左侧的两个方形为余 Cartesian 方形, 有
\begin{align}
 \delta_0\!\left(P_{\delta_1}(T,\delta_0(a))\right)
   &=P_{\delta'_2}(T,\delta'_0\delta_0(a)), \\
 \delta_1\!\left(P_{\delta_1}(T,\delta_0(a))\right)
   &=P_{\delta'_2}(T,\delta'_1\delta_0(a)).
\end{align}
因为 \(\delta'_0\delta_0=\delta'_1\delta_0\), 可得
\[
 \delta_0\!\left(P_{\delta_1}(T,\delta_0(a))\right)
 =
 \delta_1\!\left(P_{\delta_1}(T,\delta_0(a))\right);
\]
等价地,
\[
  \delta_0(\sigma_i)=\delta_1(\sigma_i)
  \qquad\text{对每个 }i.
\]

另一方面, Cayley--Hamilton 定理给出
\[
 \delta_0(a)^n-\delta_1(\sigma_1)\delta_0(a)^{n-1}
 +\cdots+(-1)^n\delta_1(\sigma_n)=0.
\]
由于 \(\delta_1(\sigma_i)=\delta_0(\sigma_i)\), 也有
\[
 \delta_0(a)^n-\delta_0(\sigma_1)\delta_0(a)^{n-1}
 +\cdots+(-1)^n\delta_0(\sigma_n)=0.
\]
因而
\[
  a^n-\sigma_1a^{n-1}+\cdots+(-1)^n\sigma_n=0,
\]
因为存在同态 \(\tau:A_1\to A_0\) 满足
\(\tau\delta_0=\operatorname{id}_{A_0}\), 从而 \(\delta_0\) 为单射.
可见 \(A_0\) 在 \(B\) 上整.
% Translation statistics: source_lines=83; output_lines=81; inline_math=22; display_math=6; equation=1; equation_star=1; align=1; footnotes=0.
% Review: complete assigned component; pending_TODO=0; all math symbols, labels, references, diagram paths and environments checked; formula prose translated.
